
\documentclass{article}
\usepackage{colortbl}
\usepackage{makecell}
\usepackage{multirow}
\usepackage{supertabular}

\begin{document}

\newcounter{utterance}

\centering \large Interaction Transcript for game `ifeval', experiment `ifeval', episode 158 with reflexion{-}llama{-}3.1{-}8b{-}instruct.
\vspace{24pt}

{ \footnotesize  \setcounter{utterance}{1}
\setlength{\tabcolsep}{0pt}
\begin{supertabular}{c@{$\;$}|p{.15\linewidth}@{}p{.15\linewidth}p{.15\linewidth}p{.15\linewidth}p{.15\linewidth}p{.15\linewidth}}
    \# & \multicolumn{2}{c}{Player} && \multicolumn{2}{c}{Game Master} \\
    \hline

    \theutterance \stepcounter{utterance}  
    & & & \multicolumn{4}{p{0.6\linewidth}}{
        \cellcolor[rgb]{0.9,0.9,0.9}{
            \makecell[{{p{\linewidth}}}]{
                \texttt{\tiny{[P1$\langle$GM]}}
                \texttt{I was hoping you could help me out with a few things. I wish to learn about how to pay my taxes. Can you summarize the process for me? Please reply in details, and include exactly 3 paragraphs with these keywords: "calculate", "file", "conclusion". Separate the paragraphs with ***.} \\
            }
        }
    }
     \\ \\

    \theutterance \stepcounter{utterance}  
    & \multicolumn{4}{p{0.6\linewidth}}{
        \cellcolor[rgb]{0.561,0.737,0.561}{
            \makecell[{{p{\linewidth}}}]{
                \texttt{\tiny{[P1$\rangle$GM]}}
                \texttt{Here's a detailed summary of the tax payment process:} \\
\\ 
\texttt{To start the tax payment process, you'll need to gather all necessary documents, including your W{-}2 forms from your employer, 1099 forms for any freelance work, and any other relevant tax{-}related documents. You'll also need to calculate your total income for the year, including any deductions and exemptions you're eligible for. This can be done using tax software or by consulting a tax professional. You'll need to determine your tax filing status, which can be single, married filing jointly, married filing separately, head of household, or qualifying widow(er). Once you have all the necessary information, you can begin calculating your tax liability using the tax tables or by using tax software.} \\
\\ 
\texttt{*** } \\
\\ 
\texttt{Once you've calculated your tax liability, you'll need to file your tax return with the IRS. You can file your return electronically or by mail, depending on your preference. If you're filing electronically, you can use tax software or the IRS's Free File program. If you're filing by mail, you'll need to complete Form 1040 and attach any necessary schedules and supporting documents. You'll also need to pay any tax due or request a refund if you've overpaid. Be sure to keep a copy of your return and any supporting documents for your records. The IRS recommends filing your return by the tax filing deadline, which is typically April 15th of each year.} \\
\\ 
\texttt{*** } \\
\\ 
\texttt{In conclusion, paying your taxes involves several steps, including gathering necessary documents, calculating your tax liability, and filing your tax return. It's essential to be accurate and thorough when completing your return to avoid any potential issues or penalties. If you're unsure about any aspect of the tax payment process, consider consulting a tax professional or using tax software to guide you through the process. Remember to keep track of your tax payments and any supporting documents, as you may need to refer to them in the future. By following these steps and staying organized, you can ensure a smooth and successful tax payment experience.} \\
            }
        }
    }
    & & \\ \\

    \theutterance \stepcounter{utterance}  
    & & & \multicolumn{2}{p{0.3\linewidth}}{
        \cellcolor[rgb]{0.9,0.9,0.9}{
            \makecell[{{p{\linewidth}}}]{
                \texttt{\tiny{[GM$|$GM]}}
                {'keywords:existence': {'keywords': ['calculate', 'file', 'conclusion']}, 'length_constraints:number_paragraphs': {'num_paragraphs': 3}}
            }
        }
    }
    & & \\ \\

    \theutterance \stepcounter{utterance}  
    & & & \multicolumn{2}{p{0.3\linewidth}}{
        \cellcolor[rgb]{0.9,0.9,0.9}{
            \makecell[{{p{\linewidth}}}]{
                \texttt{\tiny{[GM$|$GM]}}
                \texttt{game\_result = WIN} \\
            }
        }
    }
    & & \\ \\

\end{supertabular}
}

\end{document}
